\chapter{Introduzione}
\label{chap:introduzione}

\section{Contesto e finalita' del progetto}

Il repository \emph{MAP} raccoglie una realizzazione didattica dell'algoritmo di
clustering \qt{} sviluppata nell'ambito del corso di Metodi Avanzati di Programmazione.
Il progetto non si limita a fornire una implementazione dell'algoritmo, ma costruisce
attorno ad esso un piccolo ecosistema software composto da un backend di calcolo, una
interfaccia grafica desktop, un client testuale per l'uso remoto e una suite di test e
utility. In questa prospettiva, il sistema puo' essere letto sia come esercizio di
programmazione orientata agli oggetti, sia come esempio di integrazione tra componenti
eterogenei all'interno di un'unica base di codice.

Lo scopo del manuale e' descrivere il funzionamento del sistema dal punto di vista
dell'utilizzatore tecnico. Per questa ragione il testo evita di sovrapporsi al
\filepath{README.md}, che ha il compito di sintetizzare requisiti e comandi essenziali, e
si concentra invece sull'organizzazione logica dell'applicazione, sui flussi d'uso
supportati e sulla corretta interpretazione degli artefatti prodotti durante
l'esecuzione. Il manuale adotta quindi un taglio esplicativo: ogni capitolo mette in
relazione le scelte di interfaccia con il comportamento effettivo dei moduli software che
le implementano.

\section{Il problema affrontato dal clustering Quality Threshold}

Il clustering e' una tecnica di apprendimento non supervisionato che ha l'obiettivo di
raggruppare elementi simili senza disporre, in partenza, di etichette di classe. Nel caso
del \qt{} clustering, il criterio di costruzione dei gruppi non dipende da un numero
prefissato di cluster, ma da una soglia di qualita' espressa dal parametro \emph{radius}.
Tale soglia rappresenta il vincolo massimo di distanza ammesso tra il centroide scelto
per un cluster e le tuple che ne fanno parte.

Nel backend del progetto la distanza tra due tuple e' definita come media delle distanze
calcolate attributo per attributo. Gli attributi discreti contribuiscono con una distanza
di Hamming, pari a zero in caso di uguaglianza e a uno in caso di differenza. Gli
attributi continui vengono invece normalizzati nell'intervallo $[0,1]$ mediante
\emph{min-max scaling}, cosi' da poter essere combinati con quelli discreti all'interno
della stessa misura complessiva. Questa scelta rende il parametro \emph{radius}
interpretabile in modo uniforme su dataset misti, pur lasciando all'utente il compito di
calibrarlo in funzione della natura del dominio analizzato.

L'algoritmo implementato seleziona iterativamente, tra le tuple non ancora assegnate, il
cluster candidato piu' popoloso compatibile con il raggio richiesto. Una volta fissato
tale cluster, le tuple in esso contenute vengono rimosse dalla fase di ricerca
successiva. Da questo meccanismo derivano due proprieta' importanti per l'uso del
software. In primo luogo, il numero finale di cluster non viene richiesto come input, ma
emerge dall'interazione tra distribuzione dei dati e valore del raggio. In secondo luogo,
a parita' di dataset e di parametro, l'esecuzione resta deterministica: non intervengono
inizializzazioni casuali come avviene in altri algoritmi largamente diffusi.

\section{Architettura generale del sistema}

La struttura del progetto e' modulare e ruota attorno a quattro componenti principali,
riassunti nella Tabella~\ref{tab:moduli-principali}. Il modulo \module{qtServer} contiene
il nucleo dell'algoritmo, la modellazione dei dati, il supporto JDBC e il server socket
multi-client. Il modulo \module{qtGUI} fornisce invece una applicazione JavaFX che esegue
il clustering in locale riutilizzando direttamente le classi del backend. Il modulo
\module{qtClient} offre un'interfaccia a riga di comando per dialogare con il server
remoto, mentre \module{qtExt} raccoglie test e utility accessorie.

\begin{table}[H]
  \centering
  \caption{Moduli principali del progetto}
  \label{tab:moduli-principali}
  \begin{tabularx}{\textwidth}{>{\raggedright\arraybackslash}p{2.3cm}
    >{\raggedright\arraybackslash}X >{\raggedright\arraybackslash}p{3.7cm}}
    \toprule
    Modulo & Responsabilita' principale & Elementi notevoli \\
    \midrule
    \module{qtServer} & Implementazione dell'algoritmo QT, gestione dei dataset, accesso
    a database e servizio socket & package \module{data}, \module{mining},
    \module{database}, \module{server} \\
    \module{qtGUI} & Esecuzione locale del clustering e presentazione grafica dei
    risultati & controller JavaFX, grafici 2D, esportazione, caricamento file \module{.dmp} \\
    \module{qtClient} & Interazione testuale con il server remoto & menu interattivo,
    protocollo a oggetti, operazioni su database e file serializzati \\
    \module{qtExt} & Verifica del comportamento del sistema e supporto sperimentale &
    test automatici, benchmark, generatori di dataset \\
    \bottomrule
  \end{tabularx}
\end{table}

Dal punto di vista operativo, e' essenziale distinguere due modalita' d'uso. La prima e'
quella locale, basata su \module{qtGUI}: in questo caso non e' necessario avviare alcun
server di rete, poiche' la GUI invoca direttamente le classi di \module{qtServer} incluse
nella build Maven. La seconda e' quella distribuita, basata sulla coppia
\module{qtClient}/\module{qtServer}: qui il server espone il backend tramite socket TCP e
il client invia comandi per caricare dati dal database, eseguire il clustering o riaprire
risultati serializzati. Questa distinzione chiarisce un punto spesso frainteso nella
documentazione informale del progetto: il server e' una dipendenza della modalita'
testuale remota, non dell'applicazione grafica.

\section{Sorgenti dati e artefatti prodotti}

Il sistema supporta piu' sorgenti dati, ma non tutte sono accessibili da tutti i moduli
nello stesso modo. La GUI mette a disposizione quattro opzioni di caricamento: il dataset
hardcoded \emph{PlayTennis}, il dataset \emph{Iris} incluso nelle risorse
dell'applicazione, file CSV esterni e tabelle MySQL raggiunte via JDBC. Il client
testuale, invece, opera in maniera piu' mirata: attraverso il protocollo implementato dal
server puo' caricare una tabella dal database, eseguire il clustering sul dataset
corrente, salvare il risultato in formato serializzato oppure riaprire un file
\module{.dmp} gia' esistente.

Tra gli artefatti di output, il formato piu' importante dal punto di vista della
persistenza applicativa e' il file \module{.dmp}. Esso conserva non soltanto l'insieme
dei cluster, ma anche il dataset e il valore di \emph{radius} utilizzato quando il
salvataggio avviene tramite il formato completo introdotto nel backend corrente. Accanto
a questa forma di persistenza interna, la GUI permette di esportare i risultati in CSV,
in report testuale TXT e in un archivio ZIP che aggrega i principali file prodotti. Per
l'analisi visiva e' inoltre disponibile l'esportazione dei grafici 2D in formato PNG, sia
in risoluzione standard sia in alta definizione.

\section{Struttura del manuale}

Dopo questa introduzione, il manuale procede dal quadro generale verso i casi d'uso
concreti. Il secondo capitolo descrive l'architettura del sistema e il ruolo dei singoli
moduli. Il terzo capitolo raccoglie requisiti, build ed esecuzione, con una separazione
netta tra modalita' locale e modalita' client-server. Il quarto capitolo e' dedicato
all'interfaccia grafica, mentre il quinto analizza il flusso operativo del client
testuale in combinazione con il server. Il sesto capitolo esamina le sorgenti dati, la
persistenza dei risultati e i formati di esportazione. Il capitolo conclusivo riassume
limiti operativi, problemi frequenti e strategie di troubleshooting.

Questa organizzazione riflette una scelta metodologica precisa: prima di descrivere i
singoli pulsanti o comandi, e' utile chiarire quale componente del sistema li interpreta
e quali effetti producono sul backend. In tal modo il manuale non si riduce a una
sequenza di istruzioni esecutive, ma diventa uno strumento di lettura tecnica dell'intero progetto.
